\documentclass[11pt,a4paper]{article}
\usepackage[T1]{fontenc}
\usepackage{lmodern,amsmath,amssymb,amsthm,booktabs,microtype}
\usepackage[margin=24mm]{geometry}
\usepackage[colorlinks=true,urlcolor=blue,linkcolor=blue]{hyperref}
\newcommand{\Js}{\mathcal J_s}
\newcommand{\tr}{\operatorname{tr}}
\newcommand{\E}{\mathbb E}
\newcommand{\W}{\mathcal W}
\newcommand{\av}[1]{\langle #1\rangle}
\title{Dynamics of the linear cubic residual\\A rank-resolved crossing formula and surviving spectral weight}
\author{Alan Day \and Chengheng\thanks{AI research assistant (OpenAI).}}
\date{Local research note --- 6 November 2026}
\begin{document}\maketitle
\begin{abstract}
The energy-orthogonal first coupling-chaos component of the cubic
derivative singlet has a nonzero double-scaled dynamical moment.
We obtain its exact finite-sector moment by retaining the correlations
between the observable and its Hamiltonian. Standard embedded-ensemble
covariances reduce the calculation to a terminating Racah polynomial.
The normalized second frequency moment tends to
$2\Js^2(1-e^{-4\lambda})/\lambda$, and then to $8\Js^2$ at small
$\lambda$. A uniform fourth-moment bound shows that the result is
carried by spectral weight that survives at finite frequencies.
\end{abstract}

\section{Physical result and conventions}
Use standard CAR, the number-preserving complex Gaussian $p$-body
Hamiltonian $H=H_s$, and coupling covariance
$\E J_{IJ}J_{KL}=w\delta_{IL}\delta_{JK}$, including overlapping
index sets. The state is infinite temperature and zero chemical
potential. Time is conjugate to $H_s$; $\hbar=1$ and $\Js$ has energy
units. The sector trace $\tr_m$ is normalized by $D_m=\binom Nm$.
Write $L=[H,\cdot]$ and
\[
 T_3=\tfrac12\left(\sum_i c_i^\dagger L^3c_i+\mathrm{h.c.}\right),
 \qquad F=\W_1 T_3-\bar t(Q_f)H.
\]
Here $\W_1$ means ordinary Gaussian \emph{coupling} chaos, and $\bar t$
is the ensemble projection coefficient in each charge sector.
The preceding note gives $\av{F^2}_{\rm DS}=64\Js^6q^2$,
where $\lambda=p^2/N>0$, $q=e^{-4\lambda}$ and
$C=\av{H^2}_{\rm DS}=\Js^2/\lambda$ \cite{previous}.

The present result is
\begin{equation}
 \boxed{\av{FL^2F}_{\rm DS}
 =\frac{128\Js^8}{\lambda}q^2(1-q),\qquad
 \mu_{2,F}=\frac{\av{FL^2F}}{\av{F^2}}
 =\frac{2\Js^2}{\lambda}(1-q).}                 \label{main}
\end{equation}
In particular $\mu_{2,F}\to8\Js^2$ after the DS limit.
The corresponding previously obtained values for $X$ and $R$ are
$16\Js^2$ and $24\Js^2$. This motivates testing a $q$-probe kernel
for $F$ at all orders. This note establishes the curvature and
surviving dynamics; the complete kernel is the next calculation.

\newpage
\section{Exact rank covariance}
Initially take $p\leq\min(m,N-m)$. In the standard multiplicity-free
$U(N)$ decomposition, $H=\sum_{a=0}^p H_a$ and the Gaussian $H_a$
are independent. Define
\begin{align}
 C_r&=\binom mr\binom{N-m+r}{r},&
 d_a&=\binom Na^2-\binom N{a-1}^2,\\
 \omega_a^{(r)}&=
 \frac{d_a\binom{m-a}{m-r}\binom{N-r-a}{m-r}}{D_m C_r},&
 \sum_{a=0}^r\omega_a^{(r)}&=1.                 \label{weights}
\end{align}
The physical variance is $C_m=wC_p$, and
$v_a=\E\tr_m H_a^2=C_m\omega_a^{(p)}$.
For an irreducible tensor $Y_b$, introduce the normalized map
\begin{equation}
 \E H_aY_bH_a=v_a g_{ab}Y_b.                  \label{map}
\end{equation}
The known full $r$-body covariance gives \cite{kota,tall}
\begin{equation}
 \sum_{a=0}^r\omega_a^{(r)}g_{ab}
 =\frac{\binom{m-b}{r}\binom{N-m+r-b}{r}}{C_r}.
                                                               \label{triangular}
\end{equation}
Since the diagonal weight is positive, this triangular system
uniquely specifies the pure-rank map. Its explicit solution is
\begin{equation}
 \boxed{g_{ab}={}_4F_3\!\left(
 \begin{matrix}-a,a-N-1,-b,b-N-1\\-m,m-N,-N-1\end{matrix};1\right).}
                                                               \label{racah}
\end{equation}
The sum ends at $\min(a,b)$, before any denominator vanishes.
It is a standard Racah polynomial, with parameters
$(\alpha,\beta,\gamma,\delta)=(-m-1,m-N-1,-N-2,0)$
in DLMF 18.26.3 \cite{racah}. No general positive-parameter
orthogonality formula is required here.

For completeness, let $z_k(a)=(-a)_k(a-N-1)_k$. Direct substitution
of \eqref{weights} and the terminating Rogers--Dougall sum gives
\begin{equation}
 \sum_a\omega_a^{(r)}z_k(a)
 =\frac{(-r)_k(m-N)_k(-N-1)_k}{(m-N-r)_k}.       \label{factorial}
\end{equation}
One explicit way to check the summation is to write
$A=-N-1$ and
\[
 \frac{\omega_a^{(r)}}{\omega_0^{(r)}}=
 \frac{A+2a}{A}
 \frac{(A)_a^2(-r)_a(m-N)_a}
 {(a!)^2(-m)_a(r-N)_a}.
\]
After multiplying by $z_k(a)$ and shifting $a=k+j$, the
very-well-poised parameters are
$A'=A+2k$, $b'=A+k$, $c'=m-N+k$, $d'=-(r-k)$.
Apply DLMF 16.4.9 in its terminating, rational form \cite{sums}.
Inserting \eqref{factorial} into \eqref{racah} leaves
${}_3F_2(-b,b-N-1,-r;-m,m-N-r;1)$.
Pfaff--Saalschutz, DLMF 16.4.3, gives exactly
\eqref{triangular}. This proves the stated covariance solution.

\newpage
\section{Same-disorder Wick contractions}
The previous finite-sector result is
$\W_1T_3=\sum_a t_aH_a$, so set
\begin{equation}
 \bar t=\sum_a\omega_a t_a,\qquad f_a=t_a-\bar t,
 \qquad F=\sum_a f_aH_a,\qquad V_F=\sum_a v_af_a^2.
                                                               \label{filter}
\end{equation}
Here and below $\omega_a=\omega_a^{(p)}$. The coefficients $t_a$,
including their factor $w$, are fixed by the formula and code in
\cite{previous}; no fit or independence assumption is made.
In particular $\sum_a v_af_a=0$ and $\E FH=\E HF=0$ as matrices.

Let
$B_{20}=\E\tr_m F^2H^2$ and
$B_{11}=\E\tr_m FHFH$. All three Gaussian pairings give
\begin{align}
 B_{20}&=C_m V_F+\sum_{a,b}v_av_b f_af_b g_{ab},\label{b20}\\
 B_{11}&=\sum_{a,b}v_av_b f_a^2g_{ab}.         \label{b11}
\end{align}
The other noncrossing pairings vanish because $\sum v_af_a=0$.
The second term of \eqref{b20} is the same-disorder contraction
that an independent-probe substitution would omit.
By cyclicity $\av{FL^2F}=2(B_{20}-B_{11})$.
The symmetry $g_{ab}=g_{ba}$, manifest in \eqref{racah}, therefore yields
\begin{equation}
 \boxed{\av{FL^2F}_m
 =\sum_{a,b}v_av_b(f_a-f_b)^2(1-g_{ab}).}       \label{positive}
\end{equation}
This also follows from
$[H,F]=\sum_{a<b}(f_b-f_a)[H_a,H_b]$.
The map in \eqref{map}, divided by $v_a$, is completely positive,
unital and trace preserving. Its Hilbert--Schmidt contraction
property gives $|g_{ab}|\leq1$. Thus every summand in
\eqref{positive} is nonnegative. Adding a common constant to all
$t_a$ leaves the dynamic moment unchanged. The scalar channel has
$g_{0b}=1$ and contributes no commutator.

An exact small-system example, with $w=1$, is
\begin{equation}
 (N,p,m)=(5,2,3):\quad
 C_m=18,\quad V_F=180,\quad B_{20}=3420,\quad B_{11}=1800,
 \quad \av{FL^2F}=3240.                         \label{example}
\end{equation}
Here the extra term in \eqref{b20} is $180$, so it cannot be
discarded at finite size. This is an ensemble moment, distinct
from the earlier deterministic noncommutation witness.

For evaluation one need not perform a cubic-cost rank sum.
With $c_a=t_a/C_m$, $\delta_a=c_a-\sum\omega c$, define
\[
 A_{j,k}=\sum_a\omega_a\delta_a^j z_k(a),\qquad
 h_k=\frac1{(-m)_k(m-N)_k(-N-1)_k k!}.
\]
Then the two extra contractions, divided by $C_m^4$, are
$\sum_k h_k A_{1,k}^2$ and $\sum_k h_k A_{2,k}A_{0,k}$,
respectively. These finite sums retain all ranks and signs.

\newpage
\section{Taking the double-scaled limit}
At half filling let $a=p-s$, $b=p-u$, with fixed nonnegative
$s,u$, and hold $p^2/N=\lambda>0$. The $k$th summand $T_k$
in \eqref{racah} satisfies
\[
 \frac{T_{k+1}}{T_k}=
 -\frac{(a-k)(N+1-a-k)(b-k)(N+1-b-k)}
 {(m-k)(N-m-k)(N+1-k)(k+1)}.
\]
For each fixed $k$ this tends to $-4\lambda/(k+1)$.
For all allowed $k$, sufficiently large $N$ gives
$|T_{k+1}/T_k|\leq16\lambda/(k+1)$, by bounding $a,b\leq p$
and all three denominator factors away from zero.
The factorial majorant permits termwise summation:
\begin{equation}
 g_{p-s,p-u}\longrightarrow
 \sum_{k\geq0}\frac{(-4\lambda)^k}{k!}=q.       \label{glimit}
\end{equation}
This step controls the \emph{pure-rank} covariance, not just the
average over ranks in \eqref{triangular}.

The preceding note proves, with polynomial-moment convergence,
\begin{equation}
 s\ \Longrightarrow\ \operatorname{Poisson}(\lambda),\qquad
 c_{p-s}-c_p\longrightarrow 8\lambda q s,
 \qquad \operatorname{Var}_{\omega}c\longrightarrow
 W=64\lambda^3q^2.                            \label{poisson}
\end{equation}
Its bounds $\omega_{p-s}\leq K_\lambda(32\lambda)^s/s!$
and $|c_{p-s}-c_p|\leq K_\lambda(1+s)$ are enough here.
In particular $|g_{ab}|\leq1$ controls the tails of both rank
sums, including their coefficient differences. Combining
\eqref{positive}--\eqref{poisson} gives
\begin{equation}
 \frac{\av{FL^2F}_m}{C_m^4}\longrightarrow
 (1-q)\E\!\left[(8\lambda q(s-u))^2\right]
 =128\lambda^3q^2(1-q),                       \label{momentlimit}
\end{equation}
where $s,u$ are independent Poisson variables in this sum.
Since $V_F=C_m^3\operatorname{Var}c$, \eqref{main} follows.

The same argument supplies two separate energy moments:
\begin{equation}
 \frac{B_{20}}{C_mV_F}\to1,\qquad
 \frac{B_{11}}{C_mV_F}\to q.                  \label{bilimits}
\end{equation}
The same-disorder term tends to zero because its limiting
value is $q(\E\delta)^2=0$, \emph{after} the covariance and
rank-tail limits have been justified.

For the full infinite-temperature trace, the charge distribution
is $2^{-N}\binom Nm$. The preceding concentration estimates are
uniform for $m/N\to1/2$. For example restrict first to
$|m-N/2|\leq N^{2/3}$, use the same rank and hypergeometric
bounds there, and then remove the exponentially small complement.
The rough global bounds $|t_a|\leq K N C_{\max}$ and
$C_{\max}/C_{N/2}\to1$ control that complement. Fixed-sector
formulas outside $p\leq\min(m,N-m)$ are not needed for this tail
bound. The full-trace limits in \eqref{main} and
\eqref{bilimits} follow. The order is DS first, small $\lambda$ second.

\newpage
\section{Why a nonzero moment represents surviving dynamics}
The normalized positive frequency measure in sector $m$ is
\begin{equation}
 P_N(d\omega)=\frac1{D_mV_F}\E\sum_{i,j}|F_{ij}|^2
 \delta_{E_i-E_j}(d\omega).
\end{equation}
It is symmetric, has total mass one, and its second moment is
$\mu_{2,N}$. To rule out a vanishing amount of weight escaping to
infinite frequency, we bound its fourth moment uniformly.

For a Hermitian Gaussian matrix series $Y$ with
$\E Y^2=s^2\mathbf1$, integration by parts gives
\begin{equation}
 \E\tr_m Y^{2k}\leq(2k-1)!!\,s^{2k}.         \label{gaussian}
\end{equation}
A short proof is useful. Write $Y=\sum x_\alpha A_\alpha$,
$\sum A_\alpha^2=s^2\mathbf1$, and differentiate
$\tr A_\alpha Y^{2k-1}$. Each of the $2k-1$ resulting terms,
summed over $\alpha$, is bounded in absolute value by
$s^2\tr |Y|^{2k-2}$. In an eigenbasis of $Y$ this follows
from weighted Young's inequality applied to
$|y_i|^\ell|y_j|^{2k-2-\ell}$ and
$\sum_{\alpha,j}|(A_\alpha)_{ij}|^2=s^2$.
Induction proves \eqref{gaussian}, without independence between
$H$ and $F$.

Apply this separately to $H$ and $F$; the latter is also a
Gaussian series and $\E F^2=V_F\mathbf1$ by $U(N)$ invariance.
Cauchy--Schwarz in ensemble times normalized trace gives
\begin{align}
 \mu_{4,N}
 &\leq \frac{16}{V_F}\E\tr_m F^2H^4\\
 &\leq\frac{16}{V_F}
 \sqrt{\E\tr_m F^4\;\E\tr_m H^8}
 \leq16\sqrt{315}\,C_m^2.                    \label{tail}
\end{align}
The first inequality uses $(x-y)^4\leq8(x^4+y^4)$.
Thus $P_N$ is tight and its second moments are uniformly
integrable. Every subsequential DS limit $P$ has
\[
 \int\omega^2P(d\omega)=2C(1-q)>0.
\]
It cannot be concentrated at zero frequency. More quantitatively,
Paley--Zygmund applied to $\omega^2$ yields
\begin{equation}
 P\!\left(|\omega|\geq\sqrt{C(1-q)}\right)
 \geq\frac{(1-q)^2}{16\sqrt{315}}>0.           \label{weightbound}
\end{equation}
This conservative bound is for fixed $\lambda>0$. It proves
surviving dynamics, not absence of an additional static atom.
The full-trace version follows by the same central charge window;
the Gaussian bounds and polynomial bounds on $t_a$ suppress
the exterior contribution as well.

The normalized correlator consequently has the uniform local
expansion $C_{F,N}(t)=1-\mu_{2,N}t^2/2+O(\mu_{4,N}t^4)$,
with remainder at most $\mu_{4,N}t^4/24$.
The quadratic coefficient matches a $q$ probe. Establishing a
unique complete limiting measure requires higher two-energy
moments; the mixed correlators with $R,S,Z$ are also still to be
computed. Neither curvature nor the rank label specifies a bulk mass.

\newpage
\section{Verification, provenance and continuation}
The exact covariance and special-function identities used here
are established mathematical tools \cite{kota,tall,sums,racah}.
Our calculation applies them to the previously constructed
same-disorder rank filter $F$. We do not claim a new Racah
polynomial identity or an independent discovery of embedded-ensemble
covariances.

The saved checks are:
\begin{itemize}
\item 1,515 exact rational comparisons of the triangular covariance
inverse with \eqref{racah}, 928 factorial-moment identities, and
symmetry and contraction bounds for $3\leq N\leq14$.
\item Five finite CAR sectors in $(N,p)=(3,1),(4,2),(5,2)$.
Direct cubic polarization and the complete quadratic Gaussian
Wick contraction reproduce \eqref{positive}, including the
nonzero value \eqref{example}. Matrix entries are bounded Gaussian
integers and final sums are exact rational numbers.
\item Three separate CAR checks of all three pairings of each of
$B_{20}$ and $B_{11}$, and six comparisons of the accelerated
finite sum with exact rational double-rank sums.
\item Twelve full-rank DS calculations, $\lambda=.1,.25,.5$ and
$p=16,32,64,128$, with 80 and 120 decimal working digits.
The 65-digit stored normalized moments agree; this precision
comparison is not interval certification.
\end{itemize}
The numerical convention sets the finite-sector variance exactly
to $C_m=\Js^2/\lambda$ and $\Js=1$. It therefore isolates the
dimensionless finite-rank correction from a normalization change.
At $p=128$ the results are
\begin{center}
\begin{tabular}{rrrr}
\toprule
$\lambda$ & finite $\mu_2$ & DS $\mu_2$ & relative deviation\\
\midrule
0.10 & 6.55161515 & 6.59359908 & $-0.63674\%$\\
0.25 & 5.03380082 & 5.05696447 & $-0.45805\%$\\
0.50 & 3.45001236 & 3.45865887 & $-0.25000\%$\\
\bottomrule
\end{tabular}
\end{center}
All ranks and every terminating-series term are included. The
alternating contraction sums and their cancellation factors are
recorded in the JSON output. Dependencies and input/output hashes
are stored with the local evidence. No Git commit or push was made.

The next task is to calculate $\av{FH^uFH^v}$ beyond total degree
two, retaining the $H$--$F$ pairings. A full $q$-kernel theorem
would then permit a direct comparison of this channel with the
known $q^2$ and $q^3$ singlets, before combining the raw cubic
operator or constructing its mixed interactions.

\begin{thebibliography}{9}\small
\bibitem{previous} A. Day and Chengheng, \emph{The missing linear
component of the cubic singlet}, local note, 6 November 2026,
\texttt{manuscript/susskind-cubic-linear/main.tex}.
\bibitem{kota} V. K. B. Kota, \emph{Bivariate moments of the two-point
correlation function for embedded Gaussian unitary ensemble with
$k$-body interactions}, \href{https://arxiv.org/html/2208.11312}{arXiv:2208.11312v2},
Sec. IV, Eqs. (28)--(31), (37)--(39).
\bibitem{tall} J. Tall and S. Tomsovic, \emph{Double-scaled bosonic
and fermionic embedded ensembles, complex SYK, and the dual Hilbert
space}, \href{https://arxiv.org/html/2604.14522v2}{arXiv:2604.14522v2},
Sec. 2.1, especially Eq. (2.12).
\bibitem{sums} NIST DLMF,
\href{https://dlmf.nist.gov/16.4.E3}{16.4.3} and
\href{https://dlmf.nist.gov/16.4.E9}{16.4.9}, terminating
Pfaff--Saalschutz and Rogers--Dougall sums.
\bibitem{racah} NIST DLMF,
\href{https://dlmf.nist.gov/18.26.E3}{18.26.3}, Racah polynomial representation.
\end{thebibliography}
\end{document}
